\documentclass[12pt]{jlreq}
\usepackage{amsmath, amssymb}

\newcommand{\C}{\mathbb{C}}
\newcommand{\R}{\mathbb{R}}
\newcommand{\Z}{\mathbb{Z}}
\newcommand{\Tr}{\operatorname{Tr}}
\newcommand{\Hom}{\operatorname{Hom}}
\newcommand{\GL}{\operatorname{GL}}
\newcommand{\rank}{\operatorname{rank}}
\newcommand{\Id}{\operatorname{Id}}
\newcommand{\Ker}{\operatorname{Ker}}

\begin{document}

$(\Omega, \mathcal{F}, P)$ を確率空間とする。$X_n$ ($n=1,2,\dots$) は $\Omega$ 上で定義された，同じ分布を持つ独立な確率変数の列とし，
\[
  Y_n = \sum_{k=1}^n X_k,\quad Z_n = \sum_{k=1}^n 2^{-k} X_k
\]
とおく。

\begin{enumerate}
  \item $n\to\infty$ のとき $Y_n$ が概収束するための必要十分条件を $X_1$ の分布を用いて表せ。
  \item $X_1$ が可積分ならば，$n\to\infty$ のとき $Z_n$ は概収束することを示せ。
  \item $n\to\infty$ のとき $Z_n$ が概収束しないような $X_1$ の例を与え，そのことを証明せよ。
\end{enumerate}

\end{document}